\scanpage{227}
% continuation of the proof of (13)
deuxième est $\int_{a(s)}^{b(s)}\overline D\,d\varphi
=\int_{a(s)}^{b(s)}d\,\overline D\varphi
=(\overline D\varphi\unsure{]|)}(b(s))-(\overline D\varphi)(a(s))$,
\deleted{et les} et les deux expressions sont égales, grâce au fait que $a,b$ sont des sections horizontales. Notons aussi que ceci montre, si $a,b$ sont des sections \uline{horizontales}, qu'on aura
\[
 \mu(a,b)=\int_{a(s)}^{b(s)}\bigl(\sum\overline D_i\omega_i\bigr)_s
 =\varphi(b(s))-\varphi(a(s)),\tag*{(14)}
\]
où $\varphi$ est défini par la relation (10).

% sufficiency in general
Prouvons la suffisance de (7) pour la nullité de (3) en général. On va prouver ``plus généralement'' qu'on aura alors
\[
 \sum D_{i\,s}\int_{C_s}\omega_{i\,s}=0
\]
lorsque $C_s$ est un \unsure{1-cycle|1-cycles|1-chaîne} dans $X_s$, variant différentiablement avec $s$, dont le bord est ``horizontal'' en $s$, ou ce qui revient au $\overline{\mathrm m}$
\[
 \sum_i D_{i\,s}\int_{a(s)}^{b(s)}\omega_{i\,s}=0
\]
si $a,b$ sont des sections \uline{horizontales} de $X$ sur \deleted{\unsure{$S$}} \inserted{un ouvert de $S$.} \deleted{Pour} Se plaçant au voisinage d'un
% archive sheet number: 226
